% 各头独立投影；矩阵表示保留三个位置，拼接只沿特征维。
\begin{tikzpicture}[x=1mm,y=1mm,font=\figurefont\footnotesize,text=NNInk,
  nn flow/.append style={draw=NNWire,rounded corners=1.5mm},
  nn operator path/.append style={draw=NNWire,rounded corners=1.5mm},
  project/.style={nn operator,minimum width=24mm,minimum height=9mm},
  attention/.style={nn operator,minimum width=29mm,minimum height=39mm,inner sep=1pt}]
  \node[nn operator,draw=NNInput,fill=NNInput!18,minimum width=19mm,minimum height=35mm]
    (input) at (10,0) {猫\\追\\狗};
  \node[nn heading] at (10,24) {$\bm H$};
  \node[nn annotation,align=center] at (10,-25) {同一序列的\\三行表示};
  \draw[nn flow,-] (input.east)--(27,0);
  \draw[nn flow,-] (27,0)--(27,29)--(32,29);
  \draw[nn flow,-] (27,0)--(27,-29)--(32,-29);
  \fill[NNWire] (27,0) circle[radius=.45mm];
  \foreach \head/\yc in {1/29,2/-29}{
    \begin{scope}[yshift=\yc mm]
      \path[nn frame] (35,-24) rectangle (111,24);
      \node[nn heading,anchor=south west] at (36,25) {头\head};
      \node[project] (q\head) at (50,15) {$\bm W_{\mathrm q,\head}$};
      \node[project] (k\head) at (50,0) {$\bm W_{\mathrm k,\head}$};
      \node[project] (v\head) at (50,-15) {$\bm W_{\mathrm v,\head}$};
      \node[attention] (att\head) at (93,0) {缩放点积\\注意力};
      \draw[nn flow] (32,0)--(32,15)--(q\head.west);
      \draw[nn flow] (32,0)--(k\head.west);
      \draw[nn flow] (32,0)--(32,-15)--(v\head.west);
      \fill[NNWire] (32,0) circle[radius=.45mm];
      \draw[nn operator path] (q\head.east)--node[above] {$\bm Q_\head$} (78.5,15);
      \draw[nn operator path] (k\head.east)--node[above] {$\bm K_\head$} (att\head.west);
      \draw[nn operator path] (v\head.east)--node[above] {$\bm V_\head$} (78.5,-15);
    \end{scope}
  }
  \nnConcatOp[8mm]{concat}{125}{0}
  \node[nn annotation,align=center] at (140,-13) {沿特征维\\拼接};
  \draw[nn operator path] (att1.east)--node[above] {$\bm Z_1$} (125,29)--(concat.north);
  \draw[nn operator path] (att2.east)--node[below] {$\bm Z_2$} (125,-29)--(concat.south);
  \node[nn operator,minimum width=18mm] (wo) at (147,0) {$\bm W_{\mathrm o}$};
  \draw[nn operator path] (concat.east)--node[above] {$\bm C$} (wo.west);
  \node[anchor=west] (out) at (163,0) {$\bm Y$};
  \draw[nn flow,draw=NNOutput] (wo.east)--(out.west);
\end{tikzpicture}
